\section{The finite-assembly detector}\label{sec:detector}

We first isolate the representation-theoretic condition imposed by a finite
construction from the \(K(2)\)-local unit. This condition will later be
tested on the abutment of the descent spectral sequence.

Throughout the proof, \(p=3\). Fix the finite field \(k\) chosen in the
introduction. Let \(\Gamma_i\) be the one-dimensional Honda formal group of
height \(i\), in a coordinate with \(p\)-series \(T^{p^i}\), and put
\[
 D_i=\End(\Gamma_i)[1/p],\qquad
 P_i=\Aut(\Gamma_i)=\mathcal O_{D_i}^{\times},\qquad J=P_2.
\]
The Honda endomorphism-order description gives the displayed unit group
\citep[Example~2.15]{BGHS22}.
These are the unextended stabilizer groups, so their actions fix \(k\).
Choose the unramified quadratic subfield \(F\subset D_2\), and set
\[
 G_F=\mathcal O_F^\times,\quad
 \mathcal O_x=k[[x]],\quad Q_x=k((x)),\quad
 \mathbf 1_2=L_{K(2)}S.
\]
The symbol \(G_F\) denotes this unit group, not an absolute Galois group.
Its action on \(\mathcal O_x\) extends to \(Q_x\).

Let \(E_2=E_2(k)\), with \(\pi_0E_2=W(k)[[x]]\). The marked-lift
stabilizer action on its coefficients is described in
\citet[Theorem~1.1 and Equation~(1.4)]{DH95}, and its Morava \(E\)-theory
realization is supplied by \citet[Corollaries~7.6--7.7]{GH04}.
We use the cotangent convention for its degree-two periodicity line,
dual to the degree-minus-two tangent convention in those sources:
\[
 \omega=\pi_2E_2,\qquad
 \Omega=\omega/p,\qquad
 \overline\omega=\Omega\otimes_{\mathcal O_x}Q_x.
\]
Thus \(\pi_{2j}E_2=\omega^{\otimes j}\) and odd homotopy vanishes. A
negative tensor power of a line means the corresponding power of its dual.
All these lines have their semilinear \(J\)-action. Unless a relative
product is explicitly indicated, \(\wedge\) denotes the ordinary smash
product of spectra.

\begin{lemma}[Reduction modulo \(p\)]\label{lem:mod-p-nullity}
Multiplication by \(p\) is null as an \(E_2\)-module map on \(E_2/p\).
For every spectrum \(X\), it therefore annihilates the homotopy groups of
\(L_{K(2)}(E_2\wedge X)/p\). Those groups are naturally
\(\mathcal O_x\)-modules.
\end{lemma}

\begin{proof}
Regularity of \(p\) and evenness of \(E_2\) give \(\pi_1(E_2/p)=0\).
Apply \([-,E_2/p]_{E_2}\) to the cofiber sequence for multiplication by
\(p\) on \(E_2\). Restriction along \(E_2\to E_2/p\) gives an injection
\[
 [E_2/p,E_2/p]_{E_2}\hookrightarrow\pi_0(E_2/p).
\]
The class \(p\id\) maps to zero and is therefore zero. Tensoring this null
map with \(X\) and applying exact localization proves the assertion. A
chosen \(J\)-equivariant nullhomotopy is not needed: the conclusion is the
nullity of the module map and the induced action on homotopy.
\end{proof}

For a spectrum \(X\), define
\begin{equation}\label{eq:generic-test}
 \mathcal T(X)=L_{K(2)}(E_2\wedge X)/p,
 \qquad
 \mathcal V_d(X)=\pi_d\mathcal T(X)\otimes_{\mathcal O_x}Q_x.
\end{equation}
The action comes from \(E_2\), so every map of underlying spectra induces
an equivariant map on this test. The action of \(G_F\) is semilinear:
\(g(av)=g(a)g(v)\) for \(a\in Q_x\) and \(v\in\mathcal V_d(X)\).
For a character \(\chi:G_F\to k^\times\), let \(Q_x(\chi)\) denote the
semilinear line with a basis \(e_\chi\) satisfying
\(g(e_\chi)=\chi(g)e_\chi\).

We use these representations algebraically. A subrepresentation is a
\(G_F\)-stable \(Q_x\)-linear subspace; no additional continuity condition
is imposed on that subspace. Every finite-dimensional semilinear
representation has finite length, because a strict chain of subspaces has
strictly increasing dimensions. Its simple constituents are therefore
well defined up to isomorphism and multiplicity.

\begin{proposition}[Constituents forced by finite assembly]
\label{prop:constituent-test}
Suppose that \(L_{K(2)}X\) belongs to the ordinary thick subcategory
generated by \(\mathbf 1_2\) in \(K(2)\)-local spectra. Then, for every
\(d\in\mathbb Z\), the representation \(\mathcal V_d(X)\) is
finite-dimensional over \(Q_x\), and each simple constituent is isomorphic
to \(\overline\omega^{\otimes b}\) for some \(b\in\mathbb Z\).
\end{proposition}

\begin{proof}
Let \(\mathcal A\) be the full subcategory of finite-dimensional semilinear
\(G_F\)-representations over \(Q_x\) whose simple constituents have the
indicated form. It is a Serre subcategory: subobjects and quotients retain
their constituents, and the constituents of an extension are those of its
two ends. It is also closed under finite sums and retracts.

The functor \(\mathcal T\) is exact. Localization of modules from
\(\mathcal O_x\) to \(Q_x\) is exact and respects the semilinear action.
An exact triangle \(X_1\to X_2\to X_3\) consequently gives an exact sequence
\[
 \mathcal V_d(X_1)\longrightarrow\mathcal V_d(X_2)
 \longrightarrow\mathcal V_d(X_3)
 \longrightarrow\mathcal V_{d-1}(X_1)
 \longrightarrow\mathcal V_{d-1}(X_2).
\]
If the terms arising from \(X_1\) and \(X_2\) lie in \(\mathcal A\) in
every degree, then \(\mathcal V_d(X_3)\) is an extension of a subobject of
\(\mathcal V_{d-1}(X_1)\) by a quotient of \(\mathcal V_d(X_2)\), and
hence also lies in \(\mathcal A\). Rotation gives the same assertion for
any choice of two terms. The condition is therefore closed under finite
sums, shifts, cofibers, and retracts.

For the local unit, \(\mathcal T(\mathbf 1_2)\simeq E_2/p\). Indeed,
\(S\to\mathbf 1_2\) is a \(K(2)\)-equivalence, tensoring with \(E_2\)
preserves \(K(2)\)-equivalences, and \(E_2\) is \(K(2)\)-local. Even
periodicity gives
\[
 \mathcal V_{2j}(\mathbf 1_2)=\overline\omega^{\otimes j},
 \qquad \mathcal V_{2j+1}(\mathbf 1_2)=0.
\]
The same observation shows that
\(\mathcal T(X)\simeq\mathcal T(L_{K(2)}X)\) for every \(X\). The condition
therefore holds on every object of the ordinary thick subcategory generated
by \(\mathbf 1_2\), and hence on \(X\) as stated.
\end{proof}

The Proposition proves finite dimensionality as well as the restriction on
constituents. It uses neither closure under tensoring with arbitrary
modules nor an identification of dualizable objects with objects in the
ordinary thick subcategory. To contradict it, the rest of the argument
will exhibit one actual simple subquotient that is not a periodicity power.
